\section{Cross-correlations and extraction of the singular series}
\label{sec:ramanujan-correlations}

Define the finite singular series
\begin{equation}
 \mathfrak S_Q(h)=
 \sum_{q\leq Q}\frac{\mu(q)^2c_q(h)}{\varphi(q)^2}.
 \label{eq:finite-singular-series}
\end{equation}
For fixed $h\ne0$, the infinite series is absolutely convergent and equals
the prime-pair singular series
\begin{equation}
 \mathfrak S(h)=
 \sum_{q\geq1}\frac{\mu(q)^2c_q(h)}{\varphi(q)^2}.
 \label{eq:singular-series-ramanujan}
\end{equation}
Thus $\mathfrak S(h)=0$ for odd $h$, while for even $h$
\begin{equation}
 \mathfrak S(h)=
 2\prod_{p>2}\left(1-\frac1{(p-1)^2}\right)
 \prod_{\substack{p\mid h\\p>2}}\frac{p-1}{p-2}.
 \label{eq:prime-pair-singular-series}
\end{equation}

\begin{lemma}[Finite calibration and its tail]
\label{lem:finite-singular-calibration}
For fixed $h\ne0$,
\begin{equation}
 \sum_{\substack{d\leq Q\\(d,h)=1}}
 \frac{\lambda'_Q(d)}{\varphi(d)}
 =\mathfrak S_Q(h),
 \label{eq:divisor-singular-identity}
\end{equation}
and
\begin{equation}
 \mathfrak S_Q(h)=\mathfrak S(h)
 +O_h\left(\frac{(\log(2Q))^2}{Q}\right).
 \label{eq:finite-singular-tail}
\end{equation}
\end{lemma}

\begin{proof}
For the first identity, insert \eqref{eq:HB-divisor-coefficient} and
reverse the finite divisor rearrangement.  Equivalently, average
$c_q(n+h)$ over the reduced residue classes occupied by a prime $n$.  For
squarefree $q$,
\begin{equation}
 \sum_{\substack{a\ ({\rm mod}\ q)\\(a,q)=1}}c_q(a+h)
 =\mu(q)c_q(h),
 \label{eq:reduced-residue-ramanujan-average}
\end{equation}
which gives exactly \eqref{eq:divisor-singular-identity} after multiplying
by $\mu(q)/\varphi(q)^2$.

For the tail, write $g=(q,h)$.  Since $g$ ranges over the divisors of the
fixed integer $h$, the standard estimates for $1/\varphi(n)$ reduce the
absolute tail to
$O_h(\sum_{r>Q/|h|}\varphi(r)^{-2})$, which is
$O_h(Q^{-1}(\log(2Q))^2)$.  The Euler product of
\eqref{eq:singular-series-ramanujan} gives
\eqref{eq:prime-pair-singular-series}.
\end{proof}

We now compare the model with the primes.  We use the maximal
Bombieri--Vinogradov theorem in its usual smooth form: for every $A>0$
there is $B>0$ such that, if
\begin{equation}
 Q\leq \frac{X^{1/2}}{(\log X)^B},
 \label{eq:BV-Q-range}
\end{equation}
then
\begin{equation}
 \sum_{q\leq Q}\max_{(a,q)=1}\max_{y\ll_W X}
 \left|
 \sum_{\substack{n\leq y\\n\equiv a\ ({\rm mod}\ q)}}\Lambda(n)
 -\frac y{\varphi(q)}
 \right|
 \ll_{A,W}\frac{X}{(\log X)^A}.
 \label{eq:maximal-BV-used}
\end{equation}
This is a classical input \citep{IwaniecKowalski2004,MontgomeryVaughan2007}.

\begin{theorem}[Both prime--model cross-correlations]
\label{thm:prime-model-cross-correlations}
Fix $h\ne0$ and $W\in C_c^\infty((0,\infty))$.  For every $A>0$, after
choosing $B$ in \eqref{eq:BV-Q-range} sufficiently large,
\begin{align}
 \sum_n\Lambda(n)\Lambda_Q(n+h)W(n/X)
 &=\mathfrak S_Q(h)XI_0(W)
 +O_{A,h,W}\bigl(X(\log X)^{-A}\bigr),
 \label{eq:Lambda-model-cross-one}\\
 \sum_n\Lambda_Q(n)\Lambda(n+h)W(n/X)
 &=\mathfrak S_Q(h)XI_0(W)
 +O_{A,h,W}\bigl(X(\log X)^{-A}\bigr).
 \label{eq:Lambda-model-cross-two}
\end{align}
\end{theorem}

\begin{proof}
For \eqref{eq:Lambda-model-cross-one}, use
\eqref{eq:HB-divisor-expansion} to obtain progressions
$n\equiv-h\pmod d$.  If $(d,h)=1$, smooth partial summation in
\eqref{eq:maximal-BV-used} replaces the prime sum by
$XI_0(W)/\varphi(d)$.  The coefficient bound
\begin{equation}
 |\lambda'_Q(d)|\ll (\log(2Q))^2
 \label{eq:lambda-prime-cross-weight}
\end{equation}
follows from \eqref{eq:HB-divisor-pointwise}; requesting a larger exponent
in Bombieri--Vinogradov absorbs this logarithm.  The main terms give
\eqref{eq:divisor-singular-identity}.

If $(d,h)>1$, the residue class is not reduced.  Any von Mangoldt value in
that class is a power of one of the finitely many primes dividing $h$.
On the support of $W$, the total contribution is
$O_{h,W,\eps}(X^\eps)$: for each relevant prime power, sum
\eqref{eq:lambda-prime-cross-weight} only over the divisors of a fixed
integer.  This is absorbed in the asserted error.

For \eqref{eq:Lambda-model-cross-two}, the congruence is
$n+h\equiv h\pmod d$ with $d\mid n$.  The same maximal theorem and the
same non-reduced-class argument apply.  The two main terms agree because
the singular series is invariant under replacing $h$ by $-h$.
\end{proof}

Combining the cross-correlations with the model autocorrelation produces
an exact residual identity at natural scale.

\begin{theorem}[Ramanujan-calibrated correlation identity]
\label{thm:ramanujan-calibrated-correlation}
Fix $0<\delta<1/2$ and put $Q=\lfloor X^{1/2-\delta}\rfloor$.  For every
$A>0$,
\begin{align}
 &\sum_n\Lambda(n)\Lambda(n+h)W(n/X)
 \notag\\
 &\qquad={}
 \mathfrak S(h)XI_0(W)
 +\sum_nr_Q(n)r_Q(n+h)W(n/X)
 +O_{A,h,W,\delta}\bigl(X(\log X)^{-A}\bigr).
 \label{eq:Ramanujan-residual-correlation}
\end{align}
\end{theorem}

\begin{proof}
Expand $r_Q=\Lambda-\Lambda_Q$.  By
\cref{thm:prime-model-cross-correlations}, both cross terms equal
$\mathfrak S_Q(h)XI_0(W)$ up to arbitrary logarithmic savings.  By
\eqref{eq:model-autocorrelation-preliminary}, the model autocorrelation
has the same main term.  The elementary identity
\[
 \Lambda\Lambda_h
 =\Lambda_Q(\Lambda_Q)_h
 +r_Q(\Lambda_Q)_h+\Lambda_Q(r_Q)_h+r_Q(r_Q)_h
\]
therefore leaves one copy of $\mathfrak S_Q(h)XI_0(W)$.  Finally use
\eqref{eq:finite-singular-tail}; its error is a power smaller than $X$.
\end{proof}

Equation \eqref{eq:Ramanujan-residual-correlation} is not an estimate for
the residual autocorrelation.  For an even admissible shift, proving that
the last sum is $o(X)$ would already yield the Hardy--Littlewood
asymptotic.  The role of the finite model is to remove the correct local
main term before the unresolved correlation is named.
